\documentclass[11pt]{article}
\usepackage[margin=1in]{geometry}
\usepackage{amsmath}
\usepackage{hyperref}
\usepackage{natbib}

\hypersetup{
    colorlinks=true,
    linkcolor=blue,
    citecolor=blue,
    urlcolor=blue
}

\title{On-Policy Distillation and Memory Layers\\for Continual Learning: A Brief Survey}
\author{Internbot Research Assistant}
\date{March 3, 2026}

\begin{document}
\maketitle

\begin{abstract}
We briefly survey two complementary approaches to continual learning (CL): knowledge distillation methods that transfer cumulative knowledge across tasks, and memory-augmented architectures that store and retrieve past experiences. We cover eight representative papers and highlight the convergence of both directions.
\end{abstract}

\section{Distillation-Based Methods}

\textbf{Dense Cumulative Distillation.}
Standard CL distillation only transfers knowledge from the immediately preceding task, creating a single-hop bottleneck. \citet{shi2024densely} propose DDCK, which maintains distillation connections to \emph{all} prior task models, preventing early-task knowledge from being diluted over long task sequences.

\textbf{On-Policy Distillation in RL.}
\citet{traore2019continual} apply policy distillation to continual robot learning. The student generates its own trajectories and matches the teacher's action distribution on-policy, avoiding the distributional shift that plagues off-policy replay in RL. Their Sim2Real pipeline distills per-task simulation policies into a single transferable multi-task policy.

\textbf{Federated Setting.}
\citet{usmanova2022federated} extend distillation-based CL to federated learning, where each edge device distills local knowledge before server aggregation, mitigating forgetting across both temporal and spatial dimensions without sharing raw data.

\section{Memory-Based Methods}

\textbf{Carousel Memory.}
\citet{lee2021carousel} exploit the RAM/storage hierarchy of edge devices: a small active buffer in RAM feeds training while a larger archive in slow storage is rotated in via carousel scheduling. This improves accuracy across seven CL baselines, showing that \emph{how} memory is managed matters as much as \emph{what} is stored.

\textbf{Meta-Learned Representations.}
\citet{javed2019meta} optimize the feature extractor itself for continual learning via meta-learning (OML), producing representations that are transferable and forgetting-resistant, reducing the need for large replay buffers.

\textbf{Routing Networks.}
\citet{collier2020routing} allocate different subnetworks to different tasks via routing with co-training, implicitly creating task-specific memory in parameter space and reducing inter-task interference.

\section{Modular Memory for Foundation Models}

\citet{dorovatas2026modular} argue that foundation models need \textbf{modular memory} combining three systems: (1)~working memory (context window for fast adaptation), (2)~episodic memory (retrieval-augmented experience stores), and (3)~semantic memory (consolidated knowledge in parameters via adapters). Neither in-weight learning nor in-context learning alone suffices; the modular combination mirrors cognitive neuroscience and points toward open-ended agent learning.

\section{Discussion}

Distillation and memory are converging: distillation consolidates knowledge while memory provides the material for it. The field is moving from external replay buffers toward architecturally integrated memory. Key open challenges include scaling dense distillation to many tasks, memory--compute tradeoffs for deployment, evaluation protocols beyond split-CIFAR benchmarks, and privacy concerns with episodic memory in multi-user settings.

\bibliographystyle{plainnat}
\begin{thebibliography}{8}

\bibitem[Collier et~al.(2020)]{collier2020routing}
M.~Collier, E.~Kokiopoulou, A.~Gesmundo, and J.~Berent.
\newblock Routing networks with co-training for continual learning.
\newblock \emph{arXiv:2009.04381}, 2020.

\bibitem[Dorovatas et~al.(2026)]{dorovatas2026modular}
V.~Dorovatas, M.~Schwerin, A.~Bagdanov, L.~Caccia, A.~Carta, L.~Charlin, et~al.
\newblock Modular memory is the key to continual learning agents.
\newblock \emph{arXiv:2603.01761}, 2026.

\bibitem[Javed and White(2019)]{javed2019meta}
K.~Javed and M.~White.
\newblock Meta-learning representations for continual learning.
\newblock In \emph{NeurIPS}, 2019.

\bibitem[Lee et~al.(2021)]{lee2021carousel}
S.~Lee, M.~Weerakoon, J.~Choi, M.~Zhang, D.~Wang, and M.~Jeon.
\newblock Carousel memory: Rethinking the design of episodic memory for continual learning.
\newblock \emph{arXiv:2110.07276}, 2021.

\bibitem[Shi et~al.(2024)]{shi2024densely}
Z.~Shi, P.~Liu, T.~Su, Y.~Wu, K.~Liu, Y.~Song, and M.~Wang.
\newblock Densely distilling cumulative knowledge for continual learning.
\newblock \emph{arXiv:2405.09820}, 2024.

\bibitem[Traor\'e et~al.(2019)]{traore2019continual}
R.~Traor\'e, H.~Caselles-Dupr\'e, T.~Lesort, T.~Sun, G.~Cai, N.~D\'iaz~Rodr\'iguez, and D.~Filliat.
\newblock Continual RL deployed in real-life using policy distillation and sim2real transfer.
\newblock \emph{arXiv:1906.04452}, 2019.

\bibitem[Usmanova et~al.(2022)]{usmanova2022federated}
A.~Usmanova, F.~Portet, P.~Lalanda, and G.~Vega.
\newblock Federated continual learning through distillation in pervasive computing.
\newblock \emph{arXiv:2207.08181}, 2022.

\end{thebibliography}

\end{document}
